\scanpage{10}
% L1
b) Sur un \deleted{corps de car. $\unsure{=|\ne}$}\inserted{schéma $S$, où $2$ est inv.} (i.e. sur $\mathbb Z[{}^1\!/_2]$), il
% L2
y a à isom unique près, un couple $(G,\varphi)$
% L3
et un seul, où $G$ est une forme de $GP(1)_S$,
% L4
$\varphi:\unsure{\mathfrak S_4|\mathcal S_4}\longrightarrow G(S)$, avec $\rho_{0s}$ et $\rho_{1s}$ \deleted{d'ordre $3$ exactement}\inserted{non commutants}
% L5
en chaque fibre. \inserted{i.e. la repr. ne se factorise par $\mathfrak S_3$ sur aucune fibre} [En fait, \deleted{\unsure{on}} on trouve que
% L6
ladite représentation \deleted{est} \unsure{question|en question|partout} est fidèle]

% L7
B) $\unsure{\nu_0|v_0}=4,\nu_1=3$ (ou l'inverse) -- \deleted{on doit encore supp}\inserted{ici on sait que}
% L8
$\Gamma=\mathfrak S_4\times\mathbb Z_{/2}$. Le résultat général nous dit que
% L9
c'est \inserted{encore} la car. $2$ qu'il faut écarter. [On pourrait
% L10
déduire un énoncé concernant le gpe $\Gamma$ de
% L11
A), mais on préfère le traiter directement, on
% L12
trouve des renseignements suppl.) Maintenant $\underline{\tau}_0,\underline{\tau}_1,\underline{\tau}_2$
% slanted
% L12a
\margin{pas \uline{vraiment}!}
% L13
sont \deleted{d'autres}\inserted{les} \uline{gén. octaédraux} (non plus
% L14
tétraédraux) \unsure{satisfont|satisfaisant|tels que} $\rho_0^4=\rho_1^3=1$, et on
% L15
trouve

% L16
a) Sur un corps de car $2$, $\deleted{\unsure{\exists|\forall}}\forall$ hom \inserted{$\varphi:\Gamma\to G$} de
% L17
$\Gamma\simeq\mathbb Z_{/2}\times\mathfrak S_4$ dans une forme de $GP(1)_k$,
% L18-19
\deleted{\unsure{on a} $\rho_0^2=1$ ou l'ordre de $\rho_0$ est égal à $2$ ou celui $\rho_1=\mathrm{id}$}
% L19a
\inserted{ou $\rho_0$ et $\rho_1$ commutent}
% L20-22
\deleted{(i.e. : ou bien $\tau_0$ et $\tau_1$ commutent \deleted{\unsure{donc}} donc $\tau_2$ : elle se factorise par $\mathbb Z_{/2}\times\mathfrak S_3$, ou bien $\tau_1=\tau_2$ \deleted{\unsure{et}} et elle se factorise par $\mathbb Z_{/2}\times\deleted{\unsure{\mathfrak S_4}}\mathbb Z_{/2}$}
% slanted
% L22a
\inserted{i.e. elle se factorise par $\mathbb Z_{/2}\times\mathfrak S_3\simeq\mathbb D_6$}

% L23
b) Sur un schéma $S$ où $2\,1_S$ est inv,
% L24
$\exists!$ à isom unique près couple $(G,\varphi)$, avec
% L25
$\varphi:\Gamma\longrightarrow G$ forme de $GP(1)$ \inserted{tels que $\varphi_s$ ne se fac pas par $\mathbb Z_{/2}\times\mathfrak S_3\simeq\mathbb D_6$} -- et en fait
% L26
$\varphi$ est trivial sur $\mathbb Z_{/2}=\mathfrak z(\Gamma)$, et induit un
% L27
mono $\mathfrak S_4\subset G_{\unsure{\varphi|s}}(S)\qquad\forall s\in S$.
% archive sheet number: 9
